% Generated draft: requires mathematical and source review.
\documentclass[11pt]{article}
\usepackage[margin=25mm]{geometry}
\usepackage{fontspec}
\setmainfont{Latin Modern Roman}
\usepackage{amsmath,amssymb,mathtools,unicode-math}
\setmathfont{Latin Modern Math}
\usepackage{enumitem,needspace,xurl,hyperref}
\setlength{\parindent}{0pt}
\setlength{\parskip}{5pt}
\setlength{\emergencystretch}{3em}
\newcommand{\N}{\mathbb{N}}
\newcommand{\Z}{\mathbb{Z}}
\newcommand{\Q}{\mathbb{Q}}
\newcommand{\R}{\mathbb{R}}
\newcommand{\C}{\mathbb{C}}
\newcommand{\F}{\mathbb{F}}
\newcommand{\A}{\mathbb{A}}
\newcommand{\PP}{\mathbb{P}}
\newcommand{\E}{\mathbb{E}}
\newcommand{\eps}{\varepsilon}
\newcommand{\dd}{\,\mathrm{d}}
\begin{document}
\section*{UnsolvedMath ID: 800305}
\textbf{Source problem code:} catalogue entry
\subsection{Definitions and mathematical statement}
Let \(S^1=\R/\Z\), and let \(\operatorname{Diff}^{\infty}(S^1)\) denote the group of smooth circle diffeomorphisms. An action of a group \(G\) is a homomorphism \(\phi:G\to\operatorname{Diff}^{\infty}(S^1)\); it is faithful if \(\ker\phi=\{e\}\). A nonempty closed invariant set \(K\subseteq S^1\) is minimal if every orbit in \(K\) is dense in \(K\). Here an exceptional minimal set means a proper invariant Cantor set that is minimal, where a Cantor set is a nonempty compact, perfect, totally disconnected set. An action on \(S^1\) is minimal if every orbit is dense in \(S^1\).

Equip the space of actions with the topology induced by the \(C^{\infty}\) topology on the images of a finite generating set of \(G\). Equivalently, a path \(t\mapsto\phi_t\) is continuous if \(t\mapsto\phi_t(g)\) is continuous in the \(C^{\infty}\) topology for every \(g\in G\).

\textbf{Question (Navas, 2017).} For every finitely generated group \(G\) and every faithful action \(\phi:G\to\operatorname{Diff}^{\infty}(S^1)\) admitting an exceptional minimal set, does there exist a continuous path
\[
[0,1]\ni t\longmapsto\phi_t\in\operatorname{Hom}\bigl(G,\operatorname{Diff}^{\infty}(S^1)\bigr)
\]
such that \(\phi_0=\phi\), every \(\phi_t\) is faithful, every \(\phi_t\) with \(0\leq t<1\) admits an exceptional minimal set, and \(\phi_1\) is minimal?

The supplied statement does not specify orientation preservation; no such hypothesis is added here. The exceptional minimal sets may depend on \(t\), with no continuity requirement on those sets stated in the source. The supplied notation names the initial action both \(\phi_0\) and \(\phi\); these are identified above. The Cantor-set interpretation of ``exceptional minimal set'' uses the standard terminology and requires confirmation against the source.
\subsection{Short English statement}
Can every faithful smooth circle action of a finitely generated group with an exceptional minimal set be joined continuously, in the \(C^{\infty}\) topology, to a faithful minimal action, while remaining faithful throughout and retaining an exceptional minimal set at every parameter before the endpoint? The supplied registry labels this question open; this draft does not independently verify that status.
\subsection{Sources}
\begin{itemize}
\item Local UnsolvedMath registry, permanent ID \(800305\), exact source problem code \texttt{catalogue entry}, title ``Group Actions on 1-Manifolds --- Question 22.'' The record attributes the question to Navas (2017) and supplies the status ``open.''
\item Navas, \emph{Group Actions on 1-Manifolds: A List of Very Concrete Open Questions} (2017), \url{https://arxiv.org/abs/1712.06462}. The supplied source-item designation is ``Question 22 (\#21)''; the relationship between these two numbers is not explained in the record.
\item Supplied provenance: accessed 2026-10-16; extraction method \texttt{arxiv\_latex}. This access date and provenance are supplied metadata, not evidence of independent source consultation or current literature verification by this draft.
\end{itemize}
\end{document}
